\section*{Chapter \ref{ch:b2:plant-flowering} --- Reproductive Development and Flowering Control}

\begin{solution}{exo:b2:plant-flowering:1}
Vegetative \omterm{def:b2:plant-meristems:meristem}{meristem}: makes leaves with axillary buds, indefinitely
(indeterminate). \omterm{def:b2:plant-flowering:transition}{Inflorescence meristem}: makes bracts and floral
\omterm{def:b2:plant-meristems:meristem}{meristems} in their axils, indeterminate in many species. Floral
\omterm{def:b2:plant-meristems:meristem}{meristem}: makes the four whorls and stops --- determinate.
\end{solution}

\begin{solution}{exo:b2:plant-flowering:2}
Short-day plants flower when the day is shorter than a critical
length (chrysanthemum, soybean, rice); \omterm{prop:b2:plant-flowering:photoperiod}{long-day plants} when it is
longer (spinach, wheat, \emph{Arabidopsis}); day-neutral plants
\omterm{def:b2:angiosperm-reproduction:flower}{flower} on reaching a size (tomato, maize). A \omterm{prop:b2:plant-flowering:photoperiod}{short-day plant}
measures the length of the uninterrupted night.
\end{solution}

\begin{solution}{exo:b2:plant-flowering:3}
Wild type: sepal, petal, \omterm{def:b2:angiosperm-reproduction:flower}{stamen}, \omterm{def:b2:angiosperm-reproduction:flower}{carpel}. No A: \omterm{def:b2:angiosperm-reproduction:flower}{carpel}, \omterm{def:b2:angiosperm-reproduction:flower}{stamen},
\omterm{def:b2:angiosperm-reproduction:flower}{stamen}, \omterm{def:b2:angiosperm-reproduction:flower}{carpel}. No B: sepal, sepal, \omterm{def:b2:angiosperm-reproduction:flower}{carpel}, \omterm{def:b2:angiosperm-reproduction:flower}{carpel}. No C: sepal,
petal, petal, sepal, and the \omterm{def:b2:angiosperm-reproduction:flower}{flower} repeats inside.
\end{solution}

\begin{solution}{exo:b2:plant-flowering:4}
Weeks of cold and damp (winter has passed); light (the \omterm{def:b2:angiosperm-reproduction:seed}{seed} lies
near the surface); alternating temperatures (bare soil, no canopy);
leaching by rain (the rains have come); scarification (passage
through a gut or abrasion in soil); smoke and heat (a fire has
cleared the ground).
\end{solution}

\begin{solution}{exo:b2:plant-flowering:5}
16/8: night of 8 hours, below the critical 10: no. 12/12: yes.
12 light, 6 dark, flash, 6 dark: two nights of 6 hours: no.
8 light, 4 dark, 4 light, 8 dark: longest night 8 hours: no.
\end{solution}

\begin{solution}{exo:b2:plant-flowering:6}
\omterm{def:b2:plant-flowering:florigen}{Florigen} is a graft-transmissible signal made in the leaf and the
same in short- and long-day species --- the difference is in when
the leaf makes it, not what it makes. With the phloem interrupted
the graft does not induce flowering: the signal travels in the
phloem.
\end{solution}

\begin{solution}{exo:b2:plant-flowering:7}
$n = \ln(1/0.15)/0.05 = 1.90/0.05 = 38$ cold days. The spells total
$47$: the plant is vernalised on the 11th day of the third spell.
With $\alpha = 0.02$ it would need 95 days and would not be.
\end{solution}

\begin{solution}{exo:b2:plant-flowering:8}
The repressor's gene is silenced by chromatin marks that are copied
at each mitosis, so every cell of the shoot remembers the winter;
in the germ line and early embryo the marks are erased and the gene
is expressed again. Useful because a \omterm{def:b2:angiosperm-reproduction:seed}{seed} shed in summer must not
inherit its parent's permission to flower: it must wait for a
winter of its own.
\end{solution}

\begin{solution}{exo:b2:plant-flowering:9}
R: germinate (phytochrome converted to the far-red-absorbing form,
Pfr, the active one). R then FR: no (converted back to Pr).
R, FR, R: germinate. R, FR, R, FR: no. Only the last flash counts,
because each conversion is complete and the \omterm{def:b2:angiosperm-reproduction:seed}{seed} reads the final
state of the pigment.
\end{solution}

\begin{solution}{exo:b2:plant-flowering:10}
$p = 0.7$: $r = 1.28$, $1.42$, $1.40$, $-\infty$ for $g = 0.4, 0.6, 0.8,
1$: best near $g = 0.7$. $p = 0.95$: $1.99$, $2.33$, $2.55$,
$-\infty$: the optimum moves toward $g = 1$ (about $0.98$), a token
reserve only. Deserts, where good years are rare, favour heavy
\omterm{def:b2:angiosperm-reproduction:seed}{dormancy} and a large \omterm{def:b2:angiosperm-reproduction:seed}{seed} bank; meadows, where most years are good,
favour germinating nearly everything.
\end{solution}

\begin{solution}{exo:b2:plant-flowering:11}
B in all whorls: petal, petal, \omterm{def:b2:angiosperm-reproduction:flower}{stamen}, \omterm{def:b2:angiosperm-reproduction:flower}{stamen}. C in all whorls (C
represses A): \omterm{def:b2:angiosperm-reproduction:flower}{carpel}, \omterm{def:b2:angiosperm-reproduction:flower}{stamen}, \omterm{def:b2:angiosperm-reproduction:flower}{stamen}, \omterm{def:b2:angiosperm-reproduction:flower}{carpel}. No A and no B: C alone
everywhere --- \omterm{def:b2:angiosperm-reproduction:flower}{carpels} in all four whorls. C also ends the
\omterm{def:b2:plant-meristems:meristem}{meristem}; without it the centre stays indeterminate and makes \omterm{def:b2:angiosperm-reproduction:flower}{flower}
within \omterm{def:b2:angiosperm-reproduction:flower}{flower}.
\end{solution}

\begin{solution}{exo:b2:plant-flowering:12}
Day length is integrated across nights by the coincidence of light
with the clock, and induction needs several such nights; cold is
integrated as accumulating chromatin marks; the \omterm{def:b2:angiosperm-reproduction:seed}{seed} bank
integrates across years by spreading germination. An integrator
ignores a single anomalous day, a January thaw, or one bad year;
a threshold on the instantaneous value would flower the plant in a
warm week of February, or germinate every \omterm{def:b2:angiosperm-reproduction:seed}{seed} in the first rain.
\end{solution}

\begin{solution}{pb:b2:plant-flowering:1}
\textbf{1.} An 8-hour night is shorter than the critical 9.5 hours:
no flowering. To sell in August the grower must lengthen the nights
artificially with black cloth.
\textbf{2.} A 12-hour night. Opening on 15 August means induction
completed 8 weeks earlier, about 20 June; the twelve inductive
nights must start around 8 June.
\textbf{3.} Light the crop for a few minutes in the middle of each
night: a 14.5-hour night becomes two of about 7 hours, both below
the critical length. A flash suffices because phytochrome is
converted within minutes and the plant measures only the length of
uninterrupted darkness.
\textbf{4.} A flash at 21:00 leaves a dark period from
21:00 to 07:30, 10.5 hours --- longer than 9.5:
that segment induces. The break must fall so that no segment exceeds
the critical length.
\textbf{5.} One missed night breaks the run of twelve consecutive
inductive nights: the count restarts and flowering is delayed by
about as many days as had been accumulated. A cocklebur needs a
single inductive night and would already have been committed.
\textbf{6.} Phytochrome does not absorb green light, so the plant
does not see the flash; red converts it to the active form, which
reads as day; far-red immediately after converts it back, so the
night is not interrupted.
\textbf{7.} The long-day plant flowers: \omterm{def:b2:plant-flowering:florigen}{florigen} crosses the graft
union in the phloem. It identifies a mobile, graft-transmissible
signal common to both day-length classes.
\textbf{8.} $n = \ln 10/0.06 = 38.4$: 39 cold days.
\textbf{9.} 10 in November, 35 by the end of January, and the 39th
cold day falls on the fourth cold day of February: vernalised in
early February.
\textbf{10.} Six days: $f = e^{-0.36} = 0.70 > 0.10$: it stays
vegetative all summer and gives no grain --- winter wheat must be
sown in autumn.
\textbf{11.} Without the repressor it flowers without cold: it can
be sown in spring and harvested the same summer --- a spring wheat.
\textbf{12.} The silencing marks on the repressor's chromatin are
copied at each mitosis, so the apex remembers through spring; in
\omterm{def:b2:meiosis-heredity:meiosis}{meiosis} and the embryo the marks are erased, and each generation
must count its own winter.
\textbf{13.} The memory must sit in the cells that will make the
\omterm{def:b2:angiosperm-reproduction:flower}{flower} and that persist through winter --- the \omterm{def:b2:plant-meristems:meristem}{meristem} --- and it
is a state of their chromatin, not a mobile signal; the leaves that
measure day length are shed, and export a protein instead.
\textbf{14.} Class C: sepal, petal, petal, petal (the fourth whorl
becomes petals and the \omterm{def:b2:angiosperm-reproduction:flower}{flower} repeats).
\textbf{15.} No \omterm{def:b2:angiosperm-reproduction:flower}{stamens}, so no \omterm{def:b2:plant-life-cycles:heterospory}{pollen}; a few \omterm{def:b2:angiosperm-reproduction:flower}{carpels} form where C is
weakly active, so occasional \omterm{def:b2:angiosperm-reproduction:seed}{seeds} set.
\textbf{16.} Class A: \omterm{def:b2:angiosperm-reproduction:flower}{carpel}, \omterm{def:b2:angiosperm-reproduction:flower}{stamen}, \omterm{def:b2:angiosperm-reproduction:flower}{stamen}, \omterm{def:b2:angiosperm-reproduction:flower}{carpel}.
\textbf{17.} C represses A: whorl 1 becomes \omterm{def:b2:angiosperm-reproduction:flower}{carpel}, whorl 2 (B and C)
\omterm{def:b2:angiosperm-reproduction:flower}{stamen} --- \omterm{def:b2:angiosperm-reproduction:flower}{carpel}, \omterm{def:b2:angiosperm-reproduction:flower}{stamen}, \omterm{def:b2:angiosperm-reproduction:flower}{stamen}, \omterm{def:b2:angiosperm-reproduction:flower}{carpel}, the A-mutant \omterm{def:b2:angiosperm-reproduction:flower}{flower}.
\textbf{18.} The ABC genes are one family of \omterm{prop:b2:cell-differentiation:transcriptional}{transcription factors}
inherited from the common ancestor of all flowering plants; the
module was in place before roses, snapdragons and \emph{Arabidopsis}
diverged, and each has kept it.
\textbf{19.} B without A or C specifies no floral organ: a leaf-like
or bract-like organ.
\textbf{20.} $p = 0.8$: $r = 1.44$ ($g = 0.5$), $1.59$ ($0.7$), $1.55$
($0.9$).
\textbf{21.} Of the three, $g = 0.7$ (the true optimum lies near
$0.8$); at $g = 1$ the bad-year factor is zero and the lineage is
extinguished by the first bad year.
\textbf{22.} $p = 0.5$: $0.51$, $0.41$, $-0.03$: $g = 0.5$ is best,
and lower still would be better --- more \omterm{def:b2:angiosperm-reproduction:seed}{dormancy} in the desert.
\textbf{23.} The grower's \omterm{def:b2:angiosperm-reproduction:seed}{seeds} are stored, not buried, and survive
with certainty; he need not hedge within the field, but must keep a
reserve in the barn to resow after a failed year --- the same
reserve, held by him instead of by the soil.
\textbf{24.} A small \omterm{def:b2:angiosperm-reproduction:seed}{seed} carries reserves for a shoot of a
centimetre or two; germinating in the dark under soil it would die
before reaching light, so it waits for light --- which reaches it
only near the surface. Ploughing brings buried \omterm{def:b2:angiosperm-reproduction:seed}{seeds} up and
triggers a flush of weeds.
\textbf{25.} Critical night 9.5 hours, covering from about 8 June to
open on 15 August; 39 cold days; best $g \approx 0.8$ for $p = 0.8$
and about $0.5$ for $p = 0.5$.
\end{solution}
